\documentclass[11pt]{article}
\usepackage[margin=1in]{geometry}
\usepackage{amsmath,amssymb,amsthm}
\usepackage[T1]{fontenc}
\setlength{\emergencystretch}{3em}
\sloppy

\theoremstyle{plain}
\newtheorem{theorem}{Theorem}[section]
\newtheorem{lemma}[theorem]{Lemma}
\newtheorem{corollary}[theorem]{Corollary}
\theoremstyle{definition}
\newtheorem{definition}[theorem]{Definition}
\theoremstyle{remark}
\newtheorem*{remark}{Remark}

\newcommand{\Zvol}{\operatorname{Zvol}}
\newcommand{\Qvol}{\operatorname{Qvol}}
\newcommand{\Qnrm}{\operatorname{Qnrm}}
\newcommand{\supp}{\operatorname{supp}}
\newcommand{\bd}{\partial}
\newcommand{\F}{\mathbb{F}}
\newcommand{\Z}{\mathbb{Z}}
\newcommand{\Q}{\mathbb{Q}}
\newcommand{\code}[1]{\texttt{#1}}
\newcommand{\card}[1]{\lvert #1\rvert}
\newcommand{\norm}[1]{\lVert #1\rVert}

\title{Splitting, Stickerballs, and Flag Complexes for Taut Fillings\\[2pt]
\large What the Lean development \code{lean/Taut} proves}
\author{}
\date{}

\begin{document}
\maketitle

\noindent
This document states, as ordinary mathematics, \emph{exactly} what the Lean development
\code{lean/Taut} proves about taut fillings --- no strengthening inferred from prose. The
results are:
\begin{itemize}
\item[(1)] taut fillings of almost-disjoint cycles \emph{split}, and their filling volumes
\emph{add} (\S\ref{sec:split}), with a rational analogue (Corollary~\ref{cor:rational});
\item[(2)] a taut filling of a combinatorial $2$-sphere is automatically simplicial, pure, and
triangle-bounded, and forms an \emph{AnyRootedStickerball} --- a clean-shellable complex that
can be shelled starting from any tetrahedron (\S\ref{sec:main});
\item[(3)] the same filling is a \emph{flag complex} (\S\ref{sec:flag}).
\end{itemize}
Almost-disjoint splitting is the engine; the stickerball and flag-complex theorems are
inductions on the size of the filling driven by it. Typewriter names are the exact Lean
identifiers, listed with their status in \S\ref{sec:status}; everything labelled a theorem or
lemma is fully formalized with no \code{sorry} and no project axioms. Where a notion is weaker
than its classical namesake, or a classical bridge is \emph{not} formalized, this is stated
explicitly and flagged in the text as ``not Lean-proven here'' (\S\ref{sec:status}--\ref{sec:classical}).

%==========================================================================
\section{Definitions}\label{sec:defs}

Fix a linearly ordered vertex set $V$. A simplex is a finite subset $s\subseteq V$; if
$\card{s}=k+1$ we call it a $k$-simplex. A \emph{chain} is a finitely supported function
$M$ from simplices to $\Z$ (\code{Chain V}), written $M=\sum_s M(s)\,s$, with orientations
fixed by the vertex order. The boundary $\bd$ is the usual signed face operator
(\code{bdry}); $\bd\bd=0$. The \emph{norm} is $\norm{M}=\sum_s\card{M(s)}$ (\code{nrm}); a
chain is \emph{simplicial} if every coefficient lies in $\{-1,0,1\}$
(\code{SimplicialChain}), and then $\norm{M}$ is the number of simplices in its support. (For a
taut filling of a $2$-sphere, simpliciality is automatic --- Theorem~\ref{thm:simpl} --- and is
nowhere assumed.)

\begin{definition}[Filling volume; taut filling]\label{def:taut}
The \emph{filling volume} of a chain $X$ is $\Zvol(X)=\inf\{\norm{M}:\bd M=X\}$ (\code{Zvol},
literally an $\inf$ over $\mathbb N$). For a cycle ($\bd X=0$) the cone over any vertex is a filling
(\code{bdry\_cone\_of\_closed}), so the $\inf$ is attained and
\[
\Zvol(X)=\min\{\norm{M}:\bd M=X\}\qquad(\code{exists\_fill\_eq\_Zvol}).
\]
(\code{Zvol} is defined for every chain; it is used below only on cycles.)
A chain $M$ is a \emph{taut filling} of $X$ if $\bd M=X$ and $\norm{M}=\Zvol(X)$, i.e.\ $M$
is a norm-minimal filling of its own boundary (\code{IsTaut}).
\end{definition}

Tautness contributes exactly three things to the proofs, all consequences of minimality:
\begin{itemize}
\item[(T1)] \emph{Boundary control:} $\bd M=X$ holds on the nose.
\item[(T2)] \emph{Purity / no junk:} the generators of a taut filling of an $n$-cycle are
all $(n{+}1)$-simplices (\code{IsTaut.dim\_pure}), and every vertex of $M$ lies on the
boundary (\code{IsTaut.vert\_subset}, via \code{IsTaut.no\_internal\_vertex}).
\item[(T3)] \emph{No extraneous cancellation:} a taut chain contains no \emph{nontrivial}
complete cone $\operatorname{cone} x\,W$ as a subchain --- $W$ closed ($\bd W=0$), the apex $x$
absent from $W$ ($\deg_x W=0$), and some vertex $x'$ present in $W$ ($\deg_{x'} W\neq 0$)
(\code{not\_taut\_complete\_cone}); such a cone could be
refilled with strictly fewer simplices.
\end{itemize}
``Taut'' is therefore purely combinatorial $\ell^1$-minimality with these structural
corollaries; it is not a metric or curvature condition.

\begin{definition}[Combinatorial $2$-sphere]\label{def:sphere}
A set $\sigma$ of triangles is a \emph{combinatorial $2$-sphere} (\code{IsSphere2}) if it is
pure $2$-dimensional, every edge lies in exactly two triangles, every vertex link is
\emph{connected}, the $1$-skeleton is connected, and $V+F=E+2$ (that each vertex link is in fact a
cycle, i.e.\ $2$-regular, is then derived: \code{link\_two\_regular}). These properties are taken
as the definition.
\end{definition}

\begin{definition}[AlmostDisjoint]\label{def:ad}
Fix $n\ge 1$ and $A,B\subseteq V$. Closed pure $n$-cycles $X,Y$ are
\emph{AlmostDisjoint over $(A,B)$} if every generator of $X$ is an $(n{+}1)$-subset of $A$,
every generator of $Y$ is an $(n{+}1)$-subset of $B$, $\bd X=\bd Y=0$, and
\[
\card{A\cap B}\ \le\ n+1 .
\]
Mathematically: the two regions overlap in at most an $n$-simplex's worth of vertices.
This is the \emph{entire} separation hypothesis --- no metric, no topological transversality.
(Informal name; in Lean it is the hypothesis bundle of
\code{Zvol\_add\_of\_almost\_disjoint\_full} and \code{IsTaut.splits\_full}.)
\end{definition}

For the stickerball theorem we build complexes of tetrahedra one at a time. Let $\tau$ be a
set of tetrahedra ($4$-element simplices) with current boundary $B$ (a set of triangles).

\begin{definition}[Clean glue step; clean shelling]\label{def:clean}
Adding a tetrahedron $t$ to $\tau$ is a \emph{clean glue step} (\code{CleanGlueStep}) if:
\begin{itemize}
\item[(a)] $t$ meets the current boundary $B$ in one or two of its four triangular faces,
and the new boundary is $B$ updated by symmetric difference with the faces of $t$;
\item[(b)] $t\notin\tau$ and no already-interior face of $t$ is rogue (each old face of $t$
lies in a shared boundary triangle); and the three \emph{link certificates}:
\item[(c1)] each triangle of $t$ lies in at most one tetrahedron of $\tau$ (so at most two
after adding $t$);
\item[(c2)] each edge of $t$ is new to $\tau$, or shares an apex with $\tau$'s edge-link there;
\item[(c3)] each vertex of $t$ is new to $\tau$, or shares an apex with $\tau$'s vertex-link there.
\end{itemize}
A \emph{clean shelling} of $\tau$ with boundary $B$ is an ordering $t_1,\dots,t_k$ of $\tau$
in which $t_1$ is a single tetrahedron and each $t_{i+1}$ is added to $\{t_1,\dots,t_i\}$ by
a clean glue step, terminating at boundary $B$ (\code{IsCleanShelling},
\code{CleanShellFrom}).
\end{definition}

The certificates (c1)--(c3) are precisely the hypotheses needed to preserve, across one
insertion, the three conditions defining a \emph{clean $3$-complex}.

\begin{definition}[Clean $3$-complex]\label{def:clean3}
A set $\tau$ of tetrahedra is a \emph{clean $3$-complex} (\code{IsClean3Complex}) if it is
pure ($3$-dimensional), triangle-bounded (every triangle in at most two tetrahedra), and
\emph{normal}: both its edge-links and its vertex-links are connected.
\end{definition}

\begin{definition}[Stickerball; root; RootedStickerball; AnyRootedStickerball]\label{def:sb}
A set $\tau$ of tetrahedra is a \emph{stickerball} with boundary $B$ if it admits a clean
shelling ending at $B$ and $B\neq\varnothing$ (\code{IsStickerball}).
A \emph{root} of a stickerball is a tetrahedron $t\in\tau$ at which some clean shelling can
\emph{begin}. ``Rooted at $t$'' (\emph{RootedStickerball}) is not a standalone Lean predicate; it
names the existence clause for that $t$ inside \code{FreelyCleanShellable}. The stickerball is an
\emph{AnyRootedStickerball} if every tetrahedron of $\tau$ is a root: in Lean,
\code{IsAnyrootedStickerball} $=$ \code{IsStickerball} together with \code{FreelyCleanShellable},
where \code{FreelyCleanShellable} asserts a clean shelling beginning at each $t\in\tau$.
\end{definition}

\begin{remark}
``RootedStickerball'' and ``AnyRootedStickerball'' are stated in Lean-backed terms. The
clean shelling certifies normality (\S\ref{sec:classical}); we do not silently upgrade
``stickerball'' to ``shellable ball'' or ``$3$-ball'' --- the proved bridge and the missing
bridge are stated in \S\ref{sec:classical}.
\end{remark}

%==========================================================================
\section{Splitting under AlmostDisjoint}\label{sec:split}

The engine of the splitting theorem is a small-support triviality (Lemma~\ref{lem:smallsupp})
together with the no-cone property (T3) of taut chains. We also record here, once, the one
homological input --- $H_1=0$ and the cut it yields (Lemma~\ref{lem:cut}); the splitting theorem
does \emph{not} use it (its analogue of closedness is the coefficient-level
\code{bdry\_lk\_eq\_zero}, not the cut), but the sphere-cutting reductions of \S\ref{sec:main} do.

\subsection*{The homological input (used in \S\ref{sec:main})}

\begin{lemma}[$H_1=0$ and the cut]\label{lem:cut}
On a combinatorial $2$-sphere, over $\F_2$, every $1$-cycle is a boundary:
$\operatorname{im}\bd_2=\ker\bd_1$ (\code{range\_bd2\_eq\_ker\_bd1}). Consequently, if
$\gamma$ is a triple of vertices whose three edges lie in $\sigma$ but which is not a face
of $\sigma$, then the $1$-cycle of its edges bounds a set $W$ of faces of $\sigma$
(\code{exists\_cut}).
\end{lemma}

\begin{proof}
The $\F_2$ $2$-cycles are spanned by the all-faces chain (dual-graph connectivity), so
$\dim\operatorname{im}\bd_2=F-1$; the connected $1$-skeleton gives
$\dim\operatorname{im}\bd_1=V-1$, hence $\dim\ker\bd_1=E-(V-1)$, which equals $F-1$ by
$V+F=E+2$. With $\operatorname{im}\bd_2\subseteq\ker\bd_1$ this forces equality. A non-face
triple $\gamma$ has each vertex on exactly two of its edges, so its edge-chain is a
$1$-cycle and therefore a boundary $\bd_2 W$.
\end{proof}

\begin{lemma}[Small-support triviality]\label{lem:smallsupp}
A closed chain whose generators all have $n+1$ vertices, supported on a set of at most
$n+1$ vertices, is zero: with fewer than $n+1$ vertices there is no $(n{+}1)$-generator at all
(\code{eq\_zero\_of\_supp\_card\_lt}, closedness not needed), and on exactly $n+1$ vertices the
single possible generator is killed by closedness (\code{eq\_zero\_of\_closed\_supp\_card\_eq}).
\end{lemma}

\begin{proof}
On at most $n+1$ vertices there is at most one $(n{+}1)$-subset, so the chain is $c\cdot s$
for a single simplex; $\bd(c\,s)=0$ forces $c=0$.
\end{proof}

\subsection*{The split}

Fix a common vertex $p\in A\cap B$. There is a norm-non-increasing \emph{kill operator}
$K_{A,p}$ (\code{Projection.lean}, a cone projection toward $p$) with: $K_{A,p}$ fixes every
chain confined to $A$ (\code{Kmap\_eq\_self}); and $K_{A,p}$ annihilates every closed pure
$n$-cycle confined to $B$ (\code{Kmap\_eq\_zero\_of\_closed}), because its image is a closed
pure cycle on $A\cap B$, hence on $\le n+1$ vertices, hence zero by
Lemma~\ref{lem:smallsupp}. Symmetrically for $K_{B,q}$ with $q\in A\cap B$, $q\neq p$.

\begin{theorem}[Additivity]\label{thm:add}
If $X,Y$ are AlmostDisjoint over $(A,B)$ with $n\ge 1$, then
$\Zvol(X+Y)=\Zvol(X)+\Zvol(Y)$ (\code{Zvol\_add\_of\_almost\_disjoint\_full}).
\end{theorem}

\begin{proof}
$(\le)$ The sum of minimal fillings of $X$ and $Y$ is a filling of $X+Y$, so by subadditivity of
the norm, $\Zvol(X{+}Y)\le\norm{M_X}+\norm{M_Y}=\Zvol X+\Zvol Y$ (\code{Zvol\_le}, \code{nrm\_add\_le}).
$(\ge)$ Let $M$ be a minimal (hence taut) filling of $X+Y$, so $\Zvol(X{+}Y)=\norm M$; its
generators are pure $(n{+}2)$-simplices lying in $A\cup B$ (\code{IsTaut.dim\_pure},
\code{IsTaut.vert\_subset}). Each generator $t$ splits as
$(t\cap A\cap B)\sqcup(t\setminus A)\sqcup(t\setminus B)$; as $\card{A\cap B}\le n+1<n+2$,
$t$ cannot avoid both $t\setminus A$ and $t\setminus B$, so $t$ is killed by \emph{at least
one} of $K_{A,p},K_{B,q}$ (\code{Kkills\_or\_Kkills}, an inclusive ``or''). Equivalently the
two retained sets are disjoint --- a generator may be killed by both (retained by neither),
e.g.\ one with $\card{t\setminus A}=\card{t\setminus B}=2$, and that is harmless. With the
per-map bound $\norm{K_{A,p}M}+\norm{M{\restriction}_{\mathrm{kill}_A}}\le\norm M$
(\code{nrm\_Kmap\_add\_killed\_le}) and
$\norm M\le\norm{M{\restriction}_{\mathrm{kill}_A}}+\norm{M{\restriction}_{\mathrm{kill}_B}}$
(\code{nrm\_le\_killed\_add\_killed}) this gives $\norm{K_{A,p}M}+\norm{K_{B,q}M}\le\norm{M}$.
Because $K_{A,p}$ fixes the $A$-part of $\bd M$ and kills its closed $B$-part, $K_{A,p}M$ fills
$X$ and $K_{B,q}M$ fills $Y$; thus $\norm{M}\ge\Zvol X+\Zvol Y$. (Exclusivity --- each generator
killed by \emph{exactly} one --- is not used here; it holds only later, for the taut filling in
Theorem~\ref{thm:split}, once \code{no\_double\_kill} rules out double kills.)
\end{proof}

\begin{theorem}[Splitting under AlmostDisjoint]\label{thm:split}
Let $X,Y$ be AlmostDisjoint over $(A,B)$ with $n\ge 2$, and let $M$ be a taut filling of
$X+Y$. Then, writing $M_A=M\!\restriction_{\,\subseteq A}$ and
$M_B=M\!\restriction_{\,\not\subseteq A}$,
\[
M=M_A+M_B,\qquad \bd M_A=X,\qquad \bd M_B=Y,
\]
and $M_A,M_B$ are themselves taut (\code{IsTaut.splits\_full}).
\end{theorem}

\begin{proof}
Tautness enters precisely here, turning Theorem~\ref{thm:add} into rigidity. Since $M$ is
taut, $\norm{M}=\Zvol(X+Y)=\Zvol X+\Zvol Y$, while
$\norm{M}\ge\norm{K_{A,p}M}+\norm{K_{B,q}M}\ge\Zvol X+\Zvol Y$. So all inequalities are
equalities, with two consequences.

\emph{No double kill} (\code{no\_double\_kill}): if some generator were retained by neither
operator, the middle sum would be strictly smaller than $\norm M$. Hence every generator is
retained by exactly one side, and $K_{A,p}M,K_{B,q}M$ are \emph{taut} fillings of $X,Y$.

\emph{No straddling simplex} (\code{no\_extreme\_hybrid}). The three-part decomposition
leaves, besides pure generators ($\subseteq A$ or $\subseteq B$), only ``extreme hybrids'':
a single vertex $y_0$ on the far side and none in $A\cap B$. If such hybrids existed,
collect those through a fixed $y_0$; their faces opposite $y_0$ form a nonempty cycle $L$.
That $L$ is closed is \emph{not} the $\F_2$ cut: a coefficient computation shows that $\bd U$
carries no face through $y_0$ (such faces lie in neither $X$ nor $Y$), and \code{bdry\_lk\_eq\_zero}
then converts this into $\bd L=0$. So $M$ would contain the complete cone $y_0\ast L$,
contradicting (T3). This is the only place $n\ge 2$ is used:
the opposite faces need $\ge 2$ vertices for $L$ to be a genuine cycle.

Thus every generator of $M$ is $\subseteq A$ or $\subseteq B$, so the literal restrictions give
$M=M_A+M_B$. The identity $\bd M_A=X$ follows by applying $K_{A,p}$ to $\bd M_A+\bd M_B=X+Y$
(\code{Kmap\_eq\_self} on the $A$-part, \code{Kmap\_eq\_zero\_of\_closed} on the closed $B$-part);
then $\bd M_B=Y$ by cancellation. Tautness of the pieces follows from $\norm{M_A}+\norm{M_B}=\norm M$,
additivity, and $\norm{M_A}\ge\Zvol X$, $\norm{M_B}\ge\Zvol Y$.
\end{proof}

\begin{remark}[Why exactly AlmostDisjoint]
The proof uses the separation hypothesis only through Lemma~\ref{lem:smallsupp}, which needs
nothing beyond $\card{A\cap B}\le n+1$. There is no appeal to any geometric or topological
transversality; AlmostDisjointness is the exact combinatorial content that makes the cross
terms vanish.
\end{remark}

\subsection*{The rational corollary}

The integer theory transfers to rational chains. A \emph{rational chain} is a finitely supported
$M\colon\{\text{simplices}\}\to\Q$ (\code{QChain}), with the same boundary (\code{Qbdry}) and
cone (\code{Qcone}); its norm $\Qnrm M=\sum_s\card{M(s)}$ and filling volume
$\Qvol(X)=\inf\{\Qnrm M:\bd M=X\}$ are real-valued. The Lean notion of rational tautness is
\emph{order-optimal}: $M$ is \code{IsQTaut} if $\Qnrm M\le\Qnrm N$ for \emph{every} rational
chain $N$ with $\bd N=\bd M$ --- norm-minimal among co-boundary competitors. This is the exact
Lean definition; it does \emph{not} assert that the infimum $\Qvol$ is attained, so no
linear-programming or integrality theorem is needed.

\begin{corollary}[Rational additivity and splitting]\label{cor:rational}
Let $X,Y$ be rational cycles, AlmostDisjoint over $(A,B)$ (the same hypothesis bundle as
Definition~\ref{def:ad}, now over $\Q$).
\begin{itemize}
\item[\rm(i)] If $n\ge 1$ then $\Qvol(X+Y)=\Qvol X+\Qvol Y$
(\code{Qvol\_add\_of\_almost\_disjoint\_full}).
\item[\rm(ii)] If $n\ge 2$ and $M$ is an order-optimal (\code{IsQTaut}) rational filling of
$X+Y$, then $M=M_A+M_B$ with $\bd M_A=X$ and $\bd M_B=Y$, and each piece is again order-optimal
(\code{IsQTaut.splits\_full}).
\end{itemize}
\end{corollary}

\begin{proof}
Clearing denominators: every rational chain scales by some $q\in\mathbb N$ to an integer chain
(\code{exists\_nat\_smul\_integral\_three}). Part~(i) scales the cycles $X,Y$ and applies integer
additivity (Theorem~\ref{thm:add}); part~(ii) additionally scales the filling $M$ and applies
integer splitting (Theorem~\ref{thm:split}). Dividing back by $q$ gives the rational statement in
each case, order-optimality transferring because $\Qnrm$ is homogeneous under positive scaling.
\end{proof}

%==========================================================================
\section{Taut fillings are AnyRootedStickerballs}\label{sec:main}

We now prove the main theorem by induction on the number of tetrahedra, using
Theorem~\ref{thm:split} (in its $2$-sphere ``cut'' form) to produce the smaller taut
fillings. The same reduction first pays off a debt in the hypotheses: a taut filling of a
$2$-sphere is \emph{automatically} simplicial, so ``simplicial'' is a conclusion, not an
assumption.

\subsection*{Simpliciality is automatic}

\begin{theorem}[Simpliciality is automatic]\label{thm:simpl}
Let $\sigma$ be a combinatorial $2$-sphere, $X$ the $\pm1$ cycle on $\sigma$, and $M$ a taut
filling with $\bd M=X$. Then every tetrahedron coefficient of $M$ lies in $\{-1,0,1\}$ --- $M$
is simplicial (\code{taut\_filling\_is\_simplicialChain}). \emph{No simpliciality is assumed.}
\end{theorem}

This discharges the only hypothesis of the main theorem that is not already part of ``$M$ is a
taut filling of the simplicial $2$-sphere $\sigma$.'' The proof is a minimal-counterexample
deletion driven by a counting engine that produces removable tetrahedra \emph{without} knowing
the coefficients are $\pm1$ --- the point at which a naive argument would be circular.

\begin{lemma}[Disjoint eligible family, simpliciality-free]\label{lem:family}
Let $\sigma$ have no degree-$3$ vertex, and let $M$ be a taut filling whose boundary $\bd M$ is
the $\pm1$ cycle on $\sigma$ and whose generators are all tetrahedra. Then for every vertex
$v$ there are at least $\deg_v(\bd M)$ generators of $M$, pairwise disjoint in their boundary
faces, each sharing exactly two faces with $\bd M$ in proper orientation --- i.e.\ at least
$\deg_v(\bd M)$ disjoint \emph{eligible} tetrahedra
(\code{exists\_disjoint\_eligible\_family\_noS}). (Under the unit boundary an eligible
tetrahedron necessarily has coefficient $\pm1$,
\code{tet\_coeff\_eq\_pm\_one\_of\_eligible\_unitOn}; this is what makes it removable below.)
\end{lemma}

\begin{proof}
Call a triangle $\alpha\in\sigma$ a \emph{same-sign support face} of a generator $t$ if
$\alpha$ is a face of $t$ whose chain contribution carries the sign of $\bd M$ there, i.e.\
$\bd M(\alpha)\cdot M(t)\cdot\operatorname{sgn}(t,\alpha)>0$
(\code{SameSignSupportFaceTet}). Every boundary triangle is a same-sign support face of
\emph{some} generator (\code{sameSignSupportFace\_cover}); this --- strictly weaker than ``each
boundary triangle lies on a unique tetrahedron'' --- is the \emph{only} geometric input, and
it uses only that $\bd M$ is the unit cycle and that generators are tetrahedra, \emph{not} that
their coefficients are $\pm1$. Pick one such generator for each of the $\card\sigma$ boundary
triangles; since no vertex has degree $3$, each generator is picked by at most two triangles
(\code{sharedFaces\_card\_le\_two\_of\_noDegree3}).

Now count. Let $s=\card{\supp M}$ be the number of generators and $z$ the number of them with
$\card{M(t)}\ge 2$. The norm budget $s+z\le\norm M$
(\code{support\_card\_add\_nonUnitSupport\_card\_le\_nrm}; each non-unit generator costs
$\ge 2$), tautness $\norm M=\Zvol(\bd M)$, and the degree--volume inequality
$\Zvol(\bd M)+\deg_v(\bd M)\le\norm{\bd M}=\card\sigma$ (\code{Zvol\_add\_deg\_le},
\code{UnitOn.nrm\_eq}) combine to
\[
s+z+\deg_v(\bd M)\ \le\ \card\sigma .
\]
A map from the $\card\sigma$ triangles \emph{into} the $s$ generators with every fibre of size
$\le 2$ therefore has at least $z+\deg_v(\bd M)$ \emph{twofer} targets --- generators picked by
exactly two triangles --- with pairwise-disjoint assigned pairs (\code{finset\_twofer\_fibers}: a
fibred map with $s+k\le\card\sigma$ has $\ge k$ size-$2$ fibres). Discard the $\le z$ non-unit
twofers; at least $\deg_v(\bd M)$ \emph{unit} twofers remain. A unit twofer has coefficient
$\pm1$ and two same-sign boundary faces, which on a unit tetrahedron are \emph{properly
oriented} (\code{sameSignFace\_to\_proper\_of\_unit}); it is therefore eligible, and disjoint
assigned pairs make the family boundary-face disjoint.
\end{proof}

\begin{remark}
The discard is the one place coefficients re-enter, and it is essential. Counting same-sign
faces is insensitive to multiplicity --- a generator of coefficient $2$ still has each face
once --- so the twofer count runs with no simpliciality; but only a \emph{unit} twofer is
removable as an eligible tetrahedron. The norm budget pays the discard exactly: each non-unit
generator contributes the $\ge 2$ to $\norm M$ that the count demanded as surplus.
\end{remark}

\begin{proof}[Proof of Theorem~\ref{thm:simpl}]
Strong induction on $\norm M$, by the same trichotomy on $\sigma$ as in the main induction
below. If some generator $t_0$ has $\card{M(t_0)}\ge 2$
(\code{not\_simplicialChain\_iff\_exists\_natAbs\_ge\_two}), we produce a strictly smaller taut
filling of a $2$-sphere in which $t_0$ keeps its coefficient --- still non-simplicial,
contradicting the inductive hypothesis.
\begin{itemize}
\item \emph{Base} ($\card{\sigma}\le 4$): $M$ is a single tetrahedron, necessarily of
coefficient $\pm1$, so there is no defect (\code{base\_simplicialChain}).
\item \emph{Degree-$3$ vertex}: the capped cut (Lemma~\ref{lem:cut}) splits $M=M_L+M_R$ into
two strictly smaller taut fillings of the two capped spheres
(\code{taut\_splits\_for\_capped\_cut}), and $t_0$ lies on one side carrying its full
coefficient.
\item \emph{No degree-$3$ vertex}: applied at any vertex (every vertex of a $2$-sphere has
$\deg_v(\bd M)\ge 3$), Lemma~\ref{lem:family} supplies at least three eligible tetrahedra ---
in particular two distinct ones, so one is $u\neq t_0$. Remove $u$; $t_0$ persists with
coefficient $M(t_0)$
(\code{removeTet\_apply\_ne}). If $u$'s flip edge is absent, $M-u$ is a taut filling of the
flipped $2$-sphere (\code{isSphere2\_flipBoundary\_of\_eligible},
\code{isTaut\_removeTet}); if present, $M-u$ splits along that edge into two taut sub-fillings
on smaller spheres (\code{flipEdgePresent\_side\_sets\_noS},
\code{flipEdgePresent\_side\_algebra\_noS}), and $t_0$ lies on one.
\end{itemize}
In every case the smaller filling is simplicial by induction, forcing $\card{M(t_0)}\le 1$ ---
a contradiction. Hence $M$ is simplicial.
\end{proof}

\subsection*{The induction}

With simpliciality established (Theorem~\ref{thm:simpl}) we use it freely below: the Lean
lemmas \code{taut\_isPseudomanifold} and \code{theorem3\_clean} carry \code{SimplicialChain M}
as an explicit hypothesis, discharged once and for all by that theorem. First, the engine that
gives the filling its $3$-dimensional shape.

\begin{theorem}[Taut fillings are pure and triangle-bounded]\label{thm:pm}
Let $\sigma$ be a combinatorial $2$-sphere, $X$ the $\pm1$ cycle on $\sigma$
(\code{UnitOn X $\sigma$}), and $M$ a taut filling with $\bd M=X$. Then every
generator of $M$ is a tetrahedron and every triangle lies in at most two of them
(\code{taut\_isPseudomanifold}, simpliciality from Theorem~\ref{thm:simpl}; purity from
\code{IsTaut.dim\_pure}).
\end{theorem}

\begin{proof}
Purity is (T2): $X$ is $2$-dimensional, so a taut filling is $3$-dimensional.
Triangle-boundedness is a strong induction on the number of tetrahedra, by the trichotomy
used again below. \emph{Base} ($\sigma$ has $\le 4$ vertices): $M$ is a single tetrahedron.
\emph{Degree-$3$ vertex}: cut its star (Lemma~\ref{lem:cut}), apply the inductive
hypothesis to the smaller filling, re-glue. \emph{No degree-$3$ vertex}: take two disjoint
tetrahedra $e,u$ each meeting the boundary in exactly two faces
(\code{aleph\_disjoint\_eligible\_pair}), remove $e$, apply the inductive hypothesis, re-glue;
the key point is the \emph{triangle rule-out} (\code{faceCount\_removeTet\_sharedFace\_eq\_zero}):
a shared face of $e$ cannot reappear after removing $e$, since the disjoint witness $u$ would
otherwise force that triangle's incidence to $3$ in the inductively triangle-bounded filling
$M-u$.
\end{proof}

\begin{theorem}[Main]\label{thm:anyrooted}
Let $\sigma$ be a combinatorial $2$-sphere, $X$ the $\pm1$ cycle on $\sigma$, and $M$ a taut
filling with $\bd M=X$. Then the set of tetrahedra of $M$ is an
\emph{AnyRootedStickerball} with boundary $\sigma$
(\code{taut\_filling\_is\_anyrootedStickerball}). \emph{Simpliciality is not assumed}: the
endpoint derives it internally by Theorem~\ref{thm:simpl}.
\end{theorem}

By Theorem~\ref{thm:pm} the support consists of tetrahedra; $\sigma\neq\varnothing$ gives
nonempty boundary. The content is that every tetrahedron is a root:

\begin{theorem}\label{thm:freely}
Under the hypotheses of Theorem~\ref{thm:anyrooted}, for every tetrahedron $t_\star$ of $M$
there is a clean shelling of $M$ beginning with $t_\star$ (\code{theorem3\_clean}, simpliciality
from Theorem~\ref{thm:simpl}, via the induction skeleton \code{theorem3\_core\_clean} with cases
\code{base\_free\_clean}, \code{deg3\_step\_clean}, \code{prime\_step\_clean}).
\end{theorem}

\begin{proof}
Strong induction on the number $N$ of tetrahedra of $M$. The inductive invariant is exactly
the AnyRooted statement: a clean shelling beginning at \emph{any} prescribed $t_\star$. The
cases are a trichotomy on $\sigma$. In each non-base case the cut (Lemma~\ref{lem:cut})
supplies strictly smaller taut fillings of smaller spheres to which the inductive hypothesis
applies --- by Theorem~\ref{thm:split} itself in the degree-$3$ case, and by a separate
side-decomposition resting on the same cut in the prime case.

\medskip
\noindent\textbf{Base ($\sigma$ has $\le 4$ vertices).} Then $\sigma$ is a tetrahedron
boundary and $M$ is a single tetrahedron $T$; the one-term shelling starts at $T$ and is
vacuously clean.

\medskip
\noindent\textbf{Degree-$3$ vertex $v$.} The link of $v$ is a triangle $\gamma$, and the
star of $v$ is one tetrahedron $T=v\ast\gamma$. Since $\card{\sigma}\ge 5$, $\gamma$ is a
non-face triangle whose edges lie in $\sigma$, so by Lemma~\ref{lem:cut} replacing the three
triangles at $v$ by $\gamma$ yields a $2$-sphere $\sigma'$ with one fewer vertex. Restricting
$M$ to the non-star side is a \emph{direct instance of Theorem~\ref{thm:split}}, with the two
capped vertex-sets as $A,B$ (their overlap is exactly $\gamma$, so $\card{A\cap B}=n{+}1=3$):
the lemma \code{taut\_splits\_for\_capped\_cut} invokes \code{IsTaut.splits} verbatim. It
produces a \emph{taut} filling $M'$ of $\sigma'$ with strictly fewer tetrahedra, so by
induction $M'$ is an AnyRootedStickerball. Two facts let $T$ be glued cleanly: $\gamma$ is a boundary
triangle of $M'$ lying in a \emph{unique} remainder tetrahedron $t_0$ (tautness pins it,
$\text{faceCount}(\gamma)=1$; \code{degree3\_hanchor}), and $v$ is a vertex absent from $M'$
(\code{degree3\_apex\_notMem\_*}), so adjoining $T$ along the single face $\gamma$ satisfies
(a)--(c3) (the apex $v$ is new, so its link certificates are vacuous). Now assemble, given
$t_\star$:
\begin{itemize}
\item if $t_\star\neq T$: take a clean shelling of $M'$ rooted at $t_\star$ (induction) and
append $T$ last;
\item if $t_\star=T$: start with $T$, then a clean shelling of $M'$ rooted at the anchor
$t_0$ (induction), which glues to $T$ along $\gamma$.
\end{itemize}
Either way $M$ has a clean shelling rooted at $t_\star$.

\medskip
\noindent\textbf{No degree-$3$ vertex ($\card{\sigma}\ge 5$).} Pick a tetrahedron $e$ of $M$
with exactly two faces on $\sigma$ (its existence is a count: with no degree-$3$ vertex there
are more boundary triangles than tetrahedra; \code{exists\_eligibleTet}). Let $f_3,f_4$ be its
two interior faces and $f_3\cap f_4$ its \emph{flip edge}. Remove $e$.
\begin{itemize}
\item \emph{Flip edge absent.} Exposing $f_3,f_4$ is a $2$--$2$ flip of the boundary: the new
boundary is again a $2$-sphere $\sigma'$ on the same vertices
(\code{isSphere2\_flipBoundary\_of\_eligible}) and $M-e$ is a taut filling of it with fewer
tetrahedra. By induction $M-e$ is an AnyRootedStickerball; re-gluing $e$ (two shared faces)
is a clean glue step, its shared faces clean by the rule-out of Theorem~\ref{thm:pm}.
\item \emph{Flip edge present.} Then $f_3\cap f_4\in\sigma$, and $\sigma$ separates along this
edge into two $2$-spheres $\sigma_1,\sigma_2$; correspondingly $M-e$ splits along this edge into
two taut sub-fillings $M_1,M_2$ with fewer tetrahedra (\code{flipEdgePresent\_side\_sets},
\code{flipEdgePresent\_side\_algebra}). This is a decomposition of the same filter-by-side shape
as Theorem~\ref{thm:split}, but a \emph{separate} argument: it is built directly on the cut
(Lemma~\ref{lem:cut}/\code{exists\_cut}), not through \code{IsTaut.splits}.
By induction each is an AnyRootedStickerball; concatenate their clean shellings --- rooting
the second at the tetrahedron adjacent to the bridge --- and insert $e$ to bridge them
(\code{glueStep\_bridge\_left}, \code{glueStep\_bridge\_right}).
\end{itemize}
In both sub-cases, choosing the root of the smaller piece(s) appropriately (and, if
$t_\star=e$, starting with $e$) yields a clean shelling of $M$ rooted at $t_\star$.

\medskip
This proves the inductive invariant for $M$, completing the induction.
\end{proof}

\begin{proof}[Proof of Theorem~\ref{thm:anyrooted}]
Theorem~\ref{thm:freely} gives a clean shelling rooted at each tetrahedron; with nonempty
boundary this is the AnyRootedStickerball conclusion. Each re-gluing above is a clean glue
step, so the link certificates (c1)--(c3) hold at every step; the assembled shelling is
therefore a genuine clean shelling and the rooted structure is preserved through the
assembly.
\end{proof}

\begin{remark}[Where each ingredient is used]
The recursion always lands on strictly smaller taut fillings of strictly smaller spheres,
so the induction terminates. Tautness is used (i) to make every cut-restriction again taut,
so the inductive hypothesis applies (Theorem~\ref{thm:split} for the degree-$3$ cut, the
analogous side-decomposition for the prime case); (ii) through
Theorem~\ref{thm:pm} to keep every intermediate complex triangle-bounded, which the prime
rule-out needs; and (iii) to pin boundary triangles into unique tetrahedra, giving the
gluing anchors. The rooted structure is preserved because every assembly step is a
\emph{clean} glue step.
\end{remark}

%==========================================================================
\section{Taut fillings are flag complexes}\label{sec:flag}

A simplicial complex is a \emph{flag complex} when every clique of its $1$-skeleton already
spans a face: if a set of vertices has all of its edges present, the whole set is a simplex. A
taut filling inherits this.

\begin{definition}[Simplex, clique, flag complex]\label{def:flag}
For a set $\tau$ of tetrahedra, a vertex set $s$ is a \emph{simplex} of $\tau$ if $s=\varnothing$
or $s\subseteq t$ for some $t\in\tau$ (\code{SimplexOf}). It is a \emph{clique} of the
$1$-skeleton if $s\subseteq\operatorname{verts}\tau$ and every $2$-subset of $s$ is a simplex
(\code{CliqueInOneSkeleton}). Then $\tau$ is a \emph{flag complex} (\code{IsFlagComplex}) if
every clique is a simplex.
\end{definition}

\begin{theorem}[Flag complex]\label{thm:flag}
Let $\sigma$ be a combinatorial $2$-sphere, $X$ the $\pm1$ cycle on $\sigma$, and $M$ a taut
filling with $\bd M=X$. Then the tetrahedra of $M$ form a flag complex
(\code{taut\_filling\_is\_flagComplex}). \emph{Simpliciality is not assumed}: the public endpoint
derives it internally from Theorem~\ref{thm:simpl}, exactly as the stickerball endpoint does.
\end{theorem}

The proof isolates three forbidden ``taboo'' configurations and rules each out.

\begin{lemma}[No taboo $\Rightarrow$ flag]\label{lem:notaboo}
Call three vertices an \emph{empty $K_3$} if they carry all their edges yet do not span a
simplex; four vertices an \emph{empty $K_4$} likewise; and five vertices a \emph{$K_5$ of edges}
if they carry all ten edges. If $\tau$ has none of the three, it is a flag complex
(\code{NoTaboo.to\_flag}): a clique of size $0,1,2$ is a simplex outright, size $3$ or $4$ would
be an empty $K_3$ or $K_4$, and size $\ge 5$ contains a $K_5$ of edges.
\end{lemma}

\begin{proof}[Proof of Theorem~\ref{thm:flag}]
By Lemma~\ref{lem:notaboo} it suffices to rule out the three taboos for a taut filling.

\emph{No empty $K_3$ or $K_4$} (\code{no\_emptyK3K4\_of\_taut}): a strong induction on $\norm M$
by the same base / degree-$3$ / prime trichotomy as \S\ref{sec:main}. Each reduction (a
degree-$3$ peel or an eligible-tet flip) lands on a strictly smaller taut filling; an empty
$K_3$ or $K_4$ would persist into it, contradicting the inductive hypothesis.

\emph{No $K_5$ of edges} (\code{no\_k5Clique\_of\_no\_emptyK4\_taut}). This is the combinatorial
substitute for the paper's ``$S^3\not\subseteq B^3$.'' Suppose five vertices $s$ carry all ten
edges. Because there is no empty $K_4$, some $4$-subset of $s$ is an actual tetrahedron of $M$
(\code{fourSubset\_mem\_support}), so the subchain $U=M\!\restriction_{\,\subseteq s}$ is
nonzero. But with triangle-boundedness (Theorem~\ref{thm:pm}) every triangle inside $s$ is
interior, so $\bd U=0$ (\code{bdry\_filter\_subset\_eq\_zero}); and a subchain of a taut chain is
taut (\code{IsTaut.subChain}), forcing $\norm U=\Zvol(0)=0$, i.e.\ $U=0$ --- a contradiction.
\end{proof}

\begin{remark}
The $K_5$ step is where ``a taut filling has no nonzero closed subchain'' replaces the paper's
appeal to $S^3\not\subseteq B^3$; the project carries no topology, and this combinatorial
substitute is what is Lean-proven. The internal lemmas (\code{theorem4\_flag},
\code{no\_emptyK3K4\_of\_taut}) still take \code{SimplicialChain M} as data --- threaded through
their own induction --- but the public endpoint \code{taut\_filling\_is\_flagComplex} supplies it
from Theorem~\ref{thm:simpl}, so the flag-complex theorem itself carries \emph{no} simpliciality
hypothesis (its signature is \code{IsSphere2 $\sigma$ $\to$ UnitOn X $\sigma$ $\to$ $\bd$X${=}0$
$\to$ $\bd$M${=}$X $\to$ IsTaut M $\to$ IsFlagComplex M.support}).
\end{remark}

%==========================================================================
\section{Lean correspondence and status}\label{sec:status}

The development builds successfully; a search of \code{lean/Taut} finds no \code{sorry},
\code{admit}, \code{native\_decide}, \code{axiom}, or \code{unsafe}. Each theorem below
reduces to the three foundational axioms \code{propext}, \code{Classical.choice},
\code{Quot.sound} only.

\begin{center}\small
\begin{tabular}{lll}
\hline
Result here & Lean identifier & Status\\
\hline
Lem~\ref{lem:cut} ($H_1=0$, cut) & \code{range\_bd2\_eq\_ker\_bd1}, \code{exists\_cut} & fully proved\\
Thm~\ref{thm:add} (additivity) & \code{Zvol\_add\_of\_almost\_disjoint\_full} & fully proved\\
Thm~\ref{thm:split} (splitting) & \code{IsTaut.splits\_full} & fully proved\\
Thm~\ref{thm:simpl} (simpliciality automatic) & \code{taut\_filling\_is\_simplicialChain} & fully proved\\
Thm~\ref{thm:pm} (pure, triangle-bounded) & \code{taut\_isPseudomanifold} & fully proved\\
Thm~\ref{thm:freely} (rooted at each tet) & \code{theorem3\_clean} & fully proved\\
Thm~\ref{thm:anyrooted} (AnyRootedStickerball) & \code{taut\_filling\_is\_anyrootedStickerball} & fully proved\\
Thm~\ref{thm:flag} (flag complex) & \code{taut\_filling\_is\_flagComplex} & fully proved\\
Cor~\ref{cor:rational}(i) (rational additivity) & \code{Qvol\_add\_of\_almost\_disjoint\_full} & fully proved\\
Cor~\ref{cor:rational}(ii) (rational splitting) & \code{IsQTaut.splits\_full} & fully proved\\
\hline
\end{tabular}
\end{center}

\noindent\emph{Definitions weaker than their classical namesakes.}
\begin{itemize}
\item \emph{Taut} (\code{IsTaut}) is combinatorial $\ell^1$-minimality
(Definition~\ref{def:taut}), not a geometric condition.
\item \emph{Triangle-bounded} (\code{IsPseudomanifold} in the development) means only ``every
triangle in at most two tetrahedra.'' It is \emph{not} a full pseudomanifold; normality
(link-connectedness) is established separately, via the clean shelling (see
\S\ref{sec:classical}). Theorem~\ref{thm:pm} proves the triangle bound only.
\item \emph{Clean shelling} (\code{IsCleanShelling}) carries the link certificates
(c1)--(c3) as data; this is what makes it certify a clean $3$-complex, and is stronger than a
boundary-only shelling.
\end{itemize}

%==========================================================================
\section{Relation to the classical statement; remaining bridges}\label{sec:classical}

The conclusion of Theorem~\ref{thm:anyrooted} is stated in Lean-backed terms
(AnyRootedStickerball). Two questions separate this from a classical ``shellable normal
$3$-ball'' statement.

\paragraph{Proved bridge: AnyRootedStickerball $\Rightarrow$ normal $3$-pseudomanifold.}
The development proves that a stickerball is a clean $3$-complex:
\[
\code{IsStickerball.isClean3Complex}\colon\quad
\text{stickerball}\ \Longrightarrow\ \text{pure}\ \wedge\ \text{triangle-bounded}\ \wedge\
\text{Normal3},
\]
where \code{Normal3} is \emph{edge-link connected and vertex-link connected}
(\code{IsAnyrootedStickerball.isClean3Complex} for the AnyRooted form). This is proved by
iterating the single-insertion lemma \code{clean3Complex\_insert} along the clean shelling;
the three certificates (c1)--(c3) of each glue step supply exactly the triangle-count,
edge-link, and vertex-link preservation (\code{isPseudomanifold\_insert},
\code{edgeLinkConnected\_insert}, \code{vertexLinkConnected\_insert}). Thus pseudomanifoldness
(triangle bound), edge-link connectedness, vertex-link connectedness, and nonempty boundary
are all \emph{Lean-proved} consequences of the conclusion; the corresponding corollaries are
\code{taut\_filling\_is\_clean3Complex} and \code{taut\_filling\_is\_shellable}.

\paragraph{Missing bridge: stickerball $\Rightarrow$ topological/PL $3$-ball.}
The statement ``a clean-shellable normal $3$-pseudomanifold with nonempty boundary is
(PL-)homeomorphic to the $3$-ball'' is the classical shelling theorem of Danaraj--Klee /
Bj\"orner. This bridge is \textbf{not Lean-proven here}: no theorem from \code{IsStickerball} to
a topological or PL ball appears in the development (it is recorded only as a docstring remark).
Any classical reading of Theorem~\ref{thm:anyrooted} as ``$M$ triangulates a $3$-ball'' must
invoke that external theorem explicitly. In the present, fully formalized terms the conclusion is exactly:
\emph{the tetrahedra of a taut filling of a combinatorial $2$-sphere form an
AnyRootedStickerball --- a clean-shellable normal $3$-pseudomanifold, with boundary the given
sphere, that can be clean-shelled starting from any tetrahedron}.

\paragraph{Hypotheses to carry into a paper statement.}
The boundary cycle is $\pm1$ on $\sigma$ (\code{UnitOn}); the filling's simpliciality is
\emph{derived, not assumed} (Theorem~\ref{thm:simpl}, \code{taut\_filling\_is\_simplicialChain}),
so the only remaining qualifier on $M$ is tautness. Splitting (Theorem~\ref{thm:split}) needs
dimension $n\ge 2$ while additivity (Theorem~\ref{thm:add}) needs only $n\ge 1$;
AlmostDisjointness is
$\card{A\cap B}\le n+1$; the ambient $V$ is infinite; and \code{IsSphere2} includes the Euler
relation $V+F=E+2$ as part of its definition, from which $H_1=0$ (Lemma~\ref{lem:cut}) is
derived rather than assumed.

\end{document}
